\chapter{En el ojo del huracán}\label{sec:en-el-ojo-del-huracan}

Miré mis notas una última vez mientras me acercaba al atril. Sentía una extraña sensación de vacío en el corazón; con el buque insignia terrestre disparando contra nuestra capital sagrada, la evidencia de la traición humana era ahora indiscutible.

La parte más preocupante fue la facilidad con la que nuestro sistema de defensa de un billón de créditos fue destruido. Si esa nave estaba aquí para conquistarnos, las fuerzas de la Federación no podrían hacer mucho para detenerla. Pero dado que a la humanidad le gustaba presentarse como una raza dócil y pacífica, parecía más probable que intentaran salvar su imagen tras el ataque. Necesitábamos terminar nuestra asociación con la Unión Terrana mientras aún tuviéramos esa opción.

—Senadores, amigos, hoy vengo a traerles una noticia grave —hice una pausa y recorrí con la mirada el auditorio abarrotado. Todos los representantes estaban presentes, salvo por el asiento vacío del embajador xanik—. Nuestras defensas de la capital fueron bombardeadas salvajemente por una nave terrestre invasora, que está desplegando tropas mientras hablamos. Cualquiera que haya creído en una asociación a largo plazo con la humanidad puede ver ahora que sus intenciones no son benévolas.

El embajador jatari Pallum se puso de pie, molesto. —Si le das un codazo a un garat, tarde o temprano te morderá. ¿Destruiría la Federación, nos empujaría a una guerra solo para demostrar algo? Señora presidenta, recuerde que usted se lo buscó.

Dejemos que la especie militarista defienda a la humanidad; sus cerebros estaban programados para la agresión, así que no es de extrañar que se entendieran. La analogía de Pallum no era adecuada, a menos que sugiriera que los humanos no eran inteligentes. Los garats eran depredadores no sintientes del mundo jatari, domesticados para pastorear ganado.

—Si no lo supiera, diría que eso suena como una amenaza —susurré—. No se trata de demostrar nada. Se trata de una especie sedienta de sangre con armas que suponen una amenaza existencial. Si plantear preocupaciones legítimas y tratar de distanciarnos es un delito, soy culpable. La Federación nunca se unirá a un planeta de asesinos y mentirosos, no bajo mi mando.

Muchos representantes asintieron mientras hablaba. El embajador jatari pareció esforzarse por responder antes de sentarse, derrotado. Aunque la protesta de Pallum no sorprendió, sí lo hizo el silencio de la embajadora humana. Johnson estaba inclinada hacia adelante, observándome sin pestañear. Un escalofrío recorrió mi columna y respiré hondo para calmarme.

—Si creen que una especie con una historia de genocidio sistemático, guerras sangrientas y regímenes tiránicos puede cambiar, votarán para que permanezcan en la Federación —me acerqué al micrófono y bajé la voz—. Solo quería que vieran que, pese a sus mentiras y grandes discursos, no han cambiado. Siempre fue cuestión de tiempo que nos pusieran en la mira.

Me desconcertó que la embajadora Johnson se pusiera de pie y aplaudiera lentamente. Sus aplausos parecían sarcásticos, especialmente con esa sonrisa burlona.

Suspiré y golpeé el suelo con fastidio. —Embajadora Johnson, ¿tiene algo que decir? Es su derecho responder, mientras la Tierra siga siendo miembro.

—Tengo mucho que decir, pero la pregunta es si alguno de ustedes me escuchará —respondió—. En primer lugar, el lanzamiento de misiles contra su estación de defensa no fue autorizado.

 —¿Dice que su nave se volvió rebelde? No es muy tranquilizador —señalé—. ¿Qué pasará cuando la próxima tripulación decida rebelarse?

La humana me miró con enojo. —¿Provocó un asalto civil a nuestra embajada y aun así se atreve a criticar? El gobierno terrestre condena las acciones del comandante Rykov, pero es un claro caso de defensa propia. Muestre los registros de la computadora de la estación, señora presidenta. No lo hará, porque muestran que usted disparó contra la nave primero.

Murullos ansiosos recorrieron la sala. Las palabras de Johnson sembraron dudas. Consideré cuestionar quién disparó primero, pero sospeché que tendría pruebas. Si convencía a los representantes de que fue defensa propia, podría ganar votos.

—Los civiles actuaron por su cuenta. Yo simplemente...

 —¿Dijo abiertamente que los humanos no eran bienvenidos? ¿Aprobó un disturbio racial? Y mire, no niega que dispararon contra nuestra nave.

—Reconozco que el láser orbital lanzó una descarga, pero invadían nuestro espacio y se negaron a irse.

 —¿Entiende que el general de mayor rango de la Federación estaba a bordo y que la nave tenía la misión de transportarlo a la embajada? Atacó una misión de rescate diplomático. ¿Somos realmente nosotros los violentos?

Sentí que el equilibrio de la cámara se tambaleaba y maldije entre dientes. Los diplomáticos humanos tenían un don con las palabras. Me costó mantener la compostura, pero sabía que perder los estribos sería el fin.

—Sí, lo son. ¿Le importaría comentar su historia? —una sonrisa triunfante se dibujó en mi rostro. No podría defender los crímenes pasados de su especie—. ¿Sus «guerras mundiales» y la limpieza étnica? He investigado desde que lanzaron esa nanobomba.

La embajadora Johnson desvió la mirada y frunció el ceño. —Lamentamos profundamente esos años. La humanidad aprendió por las malas que la violencia no es la respuesta, y casi nos destruimos. Durante el tiempo que nos han conocido, no hemos sido esa especie. Recuerden el bien que hemos hecho, no solo el mal de nuestros años primitivos.

»¿Quién acepta más refugiados cada año? ¿Quién envía médicos a ambos lados de un conflicto? ¿Quién patrocina los derechos de los seres sintientes y redactó las leyes sobre crímenes de guerra? No sé qué más hacer para demostrar que somos pacíficos.

—Podrían dejar de construir nanobombas, por ejemplo...

—Solo las usamos para proteger a nuestros amigos. Incluso después de esta farsa, todavía los vemos como amigos. El tiempo de sus juegos terminó, señora presidenta. Tenemos problemas más grandes: una IA genocida. En lugar de sabotearnos, ¿por qué no ayuda?

Gritos de asentimiento llenaron el salón. Maldije por lo bajo. Johnson tenía respuesta para todo. ¡Los humanos eran despreciables! ¿Cómo no veían la verdad tras el ataque a la capital?

Mientras pensaba qué responder, un mensajero tujili entró en pánico. —Señora presidenta, perdóneme, pero no respondía a su holopad. Estamos bajo ataque. Cientos de acorazados descienden sobre la capital y nuestras defensas están... fuera de línea.

—¡Usted! ¡Usted lo hizo! —grité, señalando a Johnson—. Lo sabía.

La humana levantó las manos, confundida. —No somos nosotros. En serio.

 —¿Qué? ¿Otra vez naves rebeldes? ¿Cientos? —dije con desprecio.

—No son naves de la Unión Terrana. Los transpondedores las identifican como nuestras —dijo el mensajero.

Se desató el caos. Docenas de representantes gritaban a la vez. Para mí, la humanidad tenía que estar detrás. ¿Hackearon nuestras naves?

Golpeé el suelo con una pezuña pidiendo orden. —¡Silencio! ¿Han disparado? ¿Qué demandan?

—Aún no hay disparos, pero tenemos una hora para la rendición incondicional —respondió.

Miré a Johnson, evaluando su reacción. Observaba el asiento vacío del embajador xanik, como si hubiera tenido una revelación. Juraría que murmuró: «Maldita sea, Cazil».

La humana dudó y recorrió la estancia. Sus ojos se detuvieron en mí; quería oírme pedir ayuda, una disculpa servil. Pero incluso si no sabía lo de los xanik, ¿su ataque obligaría a la Unión Terrana a actuar?

—Entonces, ¿qué va a hacer, embajadora? ¿Apoderarse de nosotros? —pregunté con desprecio.

La embajadora Johnson resopló: —Tiene una forma curiosa de pedir ayuda.

Una risa amarga escapó de mi garganta. —¿Por qué nos ayudaría?

—Porque, como le he dicho, no somos lo que cree. Tengo que llamar a cierto comandante. Ese buque insignia «invasor» podría ser útil. Tiene bastantes trucos bajo la manga.

Tras dudar, asentí a regañadientes. Un gesto humano, una concesión. No confiaba en ellos, pero no había opción. Quedaba ver si la Unión Terrana se enfrentaría a su mayor aliado, pero no iba a contener el aliento.